% Section "imitation" (draft of the imitation synthesis, research/tracks/imitation/notes-final.md §2; placement: after
% sec:limits, before sec:computable). Sources: research/tracks/imitation/robust/notes-final.md (Lemma 2.1, Thm 2.2,
% Props 3.1, 3.3, Setting 4.1, Lemmas 4.4-4.6, Thm 4.7(U),(vii), Cor 4.9, Props 4.10(a2),(c),(d), 4.11, Props 5.2, 5.3)
% with referee-A.md (A1, R1) and referee-B.md (B1-B3); imitation/notes-final.md §2 (B_beta conditioned on (h_0, mu) only;
% the proof of Thm 4.7(U) needs no symmetry of information). Local notation: Sigma (coin set; B is the truncation level of
% sec:teacher), rho (coin law; not a structure function), beta (coupling), W (cells), bar mu, P^*.
\section{Imitation: a sampler whose coins the environment drives}\label{sec:imit}

Imitation is one motivation for the first reason: a learner trained on a person's answers should return a copy of the person, a sampler with its own randomness. A sampler with some of its sampled bits stored is rejected per query (\cref{thm:recovery:main}). A system that predicts the answers by inferring, from the data seen so far, which person it faces is a single hypothesis only in a class whose hypotheses carry information across queries; there a Bayes mixture over persons that gives the person weight $w$ loses at most $\log(1/w)$ bits of fit to the person's own program on every data set, private coins or not, and so beats it at every $n$ once $\lambda$ times the length it saves exceeds $\log(1/w)$ (as in \cref{thm:split:shared}). The private-randomness class excludes such hypotheses by construction (\cref{cor:absorb:nouniversal}), not because of anything in the data, and a person with memory is not a product sampler. Here we treat hypotheses that predict the person's answers \emph{better} than the person's own kernel does, because the person's coins depend on the environment and they model the environment. Proofs are in \cref{app:imit}.

\begin{remark}[Private coins]\status{proved}\src{imitation/robust notes-final Lemma 2.1, Thm 2.2(2),(3); check r1; referee rb1}\label{rem:imit:private}
If, given its input, the person's answer is independent of everything a hypothesis may use (the earlier data and any variables of the environment, earlier or later), the argument of \cref{lem:split:ville} with this larger filtration shows that no predictor, however it models the environment, gains over the person more than its length plus a Ville deficiency $\delta^*$ with $\Prob[\delta^*\ge j]\le2^{-j}$ (\cref{app:imit:private}). For product hypotheses that also see an environment variable in a countable set, with i.i.d.\ pairs, none then predicts better in population, and \cref{thm:recovery:main} returns the person. This is the degenerate case: a hypothesis can only be lucky on a sample, at one bit of length per bit of luck.
\end{remark}

\begin{definition}[Environment-driven coins]\src{imitation/robust notes-final Setting 4.1 (B$_\beta$ conditioned on $(h_0,\mu)$ only, imitation/notes-final \S2.2); referee-B B1}\label{def:imit:setting}
Let $Y$ and the coin set $\Sigma$ be finite, $m\ge1$, and $W:=\{0,1\}^m$, the \emph{cells}.
(a) A \emph{channel} $\mu=(\mu(\cdot\mid w))_{w\in W}$ of finitely supported laws on inputs with rational masses gives the input law $\bar\mu:=2^{-m}\sum_w\mu(\cdot\mid w)$, with support $\supp\mu$; in the \emph{feature case} every input in $\supp\mu(\cdot\mid w)$ has $w$ as its first $m$ bits.
(b) The \emph{person} has a coin law $\rho$ on $\Sigma$, every $\rho(b)2^m$ a positive integer, and a response function $h:\supp\mu\times\Sigma\to Y$, injective in $b$; its kernel is $\nu_H(y\mid x):=\rho(b)$ if $y=h(x,b)$, and $0$ otherwise, and its \emph{model} is $h_0:=(h,\rho)$. A \emph{copy} is a program of $\Hcl{\times}$ whose kernel equals $\nu_H$ on $\supp\mu$; $\mathcal E$ is the set of copies, $\ell_0:=\min_{\mathcal E}\ell$, and $\pi_H$ is a shortest copy.
(c) A \emph{coupling} $\beta:W\to\Sigma$ of exact type $\rho$, $|\beta^{-1}(b)|=\rho(b)2^m$, drives the coins: at each query, independently, $w_i$ is uniform on $W$, $x_i\sim\mu(\cdot\mid w_i)$ and $y_i:=h(x_i,\beta(w_i))$. The learner sees $(x_i,y_i)$, i.i.d.\ from $\bar\mu\otimes P^*$ with $P^*(y\mid x):=\sum_w\mu(x\mid w)\ind[y=h(x,\beta(w))]/\sum_w\mu(x\mid w)$.
(d) $R(\pi):=\Ex[-\log\nu_\pi(y_i\mid x_i)]$, $J_c:=c\,\ell+R$ and $\Phi(c):=\inf_{\Hcl{\times}}J_c$; every copy has $R=H(\rho)$. $L_\beta:=\log\big(2^m!/\prod_b(\rho(b)2^m)!\big)$, $B_\beta:=K(\beta\mid h_0,\mu)$, $\delta_\beta:=L_\beta-B_\beta$, and $I_\beta(\pi):=B_\beta-K(\beta\mid h_0,\pi,\mu)$, the information of $\pi$ about the coupling. For a probability $P$ and a sub-probability $Q$, $D_{1/2}(P\Vert Q):=-2\log(1-\Hel(P,Q))$.
\end{definition}

So the person draws a coin at every query, but the coin it uses is a fixed function of the cell, which the learner sees in the input (the feature case) or only through the channel (a latent shared cause). The data identify $P^*$ (\cref{thm:recovery:main}); when the input carries information about the cell, hypotheses that model $\beta$ predict better than the person, by up to $H(\rho)$ bits per sample in the feature case, where $P^*$ is deterministic. The teacher of \cref{def:teacher:instance} is the feature case with $h(x,b)=f_a(x)\oplus b$, $\rho=\Bern(p)$ and $\beta=\ind_S$: systematic errors are environment-driven coins. As for error rules, $L_\beta\ge2^mH(\rho)-|\Sigma|\log(2^m+1)$ \citep{cover2006elements}, and $\Prob[\delta_\beta>j]<2^{-j}$ for a coupling drawn uniformly from its type class.

\begin{theorem}[Robust near-imitation]\status{proved}\src{imitation/robust notes-final Thm 4.7(U),(vii), Lemmas 4.4--4.6; referee-A A1; imitation/notes-final \S2.2 (re-derived); checks f7, s1(B),(C)}\label{thm:imit:main}
In the setting of \cref{def:imit:setting} let $c=a2^{-m}$ with $a>1$, $\epsilon\ge0$, and $G:=\ell_0+\delta_\beta+|\Sigma|\log(2^m+1)+\kappa$. Every $\pi\in\Hcl{\times}$ with $J_c(\pi)\le\Phi(c)+\epsilon$, whatever its internal structure, satisfies:
\begin{enumerate}[label=(\roman*),itemsep=0pt,topsep=2pt]
\item $\ell(\pi)\le\ell_0+(G+2^m\epsilon)/(a-1)$;
\item $-(aG+2^m\epsilon)/((a-1)2^m)\le R(\pi)-H(\rho)\le c\,\ell_0+\epsilon$: it predicts neither much better nor much worse than the person;
\item $I_\beta(\pi)\le\ell_0+(G+2^m\epsilon)/(a-1)+\kappa$;
\item with $E:=2\log\big(2^{m-1}(2H(\rho)+c\,\ell_0+\epsilon)+1\big)\le2m+2\log(2H(\rho)+c\,\ell_0+\epsilon+1)$,
\[
\Ex_{x\sim\bar\mu}D_{1/2}\big(\nu_H(\cdot\mid x)\,\Vert\,\nu_\pi(\cdot\mid x)\big)\ \le\ c\,\ell_0+\frac{2a}{a-1}\cdot\frac{G+E+\kappa}{2^m}+\frac{a+1}{a-1}\,\epsilon,
\]
so its Hellinger gap from the person, $\Ex_{x\sim\bar\mu}\Hel(\nu_H(\cdot\mid x),\nu_\pi(\cdot\mid x))$, is at most $\frac{\ln2}2$ times the right side.
\end{enumerate}
Moreover $J_c(\pi_H)\le\Phi(c)+G/2^m+c\,\ell_0$. In the feature case, if $|\Sigma|\ge2$ and the values of $\rho$ are distinct, a program of length at most $\ell_0+\lceil L_\beta\rceil+a_\rho+\kappa$, $a_\rho:=2\log(m+1)+2\log(|\Sigma|+1)+|\Sigma|(m+1)$, has kernel $P^*$ and risk $0$, and so beats every copy for $c<H(\rho)/(\lceil L_\beta\rceil+a_\rho+\kappa)$, a threshold at least $2^{-m}(1+(a_\rho+1+\kappa)/(2^mH(\rho)))^{-1}$.
\end{theorem}

\emph{Proof idea.} A program is a code for the coupling: $K(\beta\mid h_0,\pi,\mu)\le2^mR(\pi)+\kappa$ (\cref{lem:imit:coding}, which extends \cref{lem:teacher:coding}). With the size of the type class this gives the whole class a right line of slope $2^{-m}$ and offset $G/2^m$ at $(\ell_0,H(\rho))$, and \cref{lem:kink:main} with $\mathcal N=\emptyset$ gives (i), (ii) and the bound on $\Phi$. For (iv) a closeness lemma (\cref{lem:imit:closeness}) bounds the divergence by that of per-cell geometric means of the predictions, which define a semimeasure on couplings whose value at $\beta$ is too small unless the divergence is small or the program predicts better or knows more about $\beta$.

The best predictor has $\Ex_xD_{1/2}(\nu_H\Vert P^*)=H(\rho)$ in the feature case. So when the coupling is incompressible ($\delta_\beta\ll2^mH(\rho)$) and the person is much simpler than the coupling ($\ell_0+|\Sigma|m\ll2^mH(\rho)$, about the coupling's $L_\beta$ bits), for every $c$ between the reuse factor $2^{-m}$ and about $H(\rho)/\ell_0$, with $a$ bounded away from $1$, every near-optimal hypothesis is close to the person in behaviour, on the scale $H(\rho)$, and holds little of the coupling, although hypotheses that model the environment predict better. Below $2^{-m}$ the program that models the coupling wins; above about $H(\rho)/\ell_0$ the bound exceeds $H(\rho)$, and above about $(\log|Y|-H(\rho))/\ell_0$ a simpler, worse program can win, as above $c_{\rm hi}$ in \cref{sec:teacher}. The theorem gives near-copies, not copies: in $\Hcl{\times}$ nothing excludes a nearby coin law as the minimiser (cf.\ \cref{rem:teacher:selected}); and the hypotheses that predict better gain only by storing an incompressible table, while compact models of how the environment drives the person are the negative case of \cref{prop:imit:limits}(c). By \cref{prop:recovery:beyond}(c) with the sampler $P^*$, almost surely every $cn$-MAP satisfies (i)--(iv) with $\epsilon=0$ for all large $n$. For example, with $m=20$, $\epsilon=\delta_\beta=0$ and the machine constant set to $0$ (an illustration), a direct computation (code in the accompanying repository) gives at $c=2^{-19}$ the bound $6.2\cdot10^{-4}$ bits for a person who answers with a fair coin ($h(x,b)=b$, $\ell_0\le a_\rho<54$), against $H(\rho)=1$, and $9.5\cdot10^{-4}$ for the parity person of \cref{sec:teacher} ($\ell_0<120$), against $H(\rho)=0.0112$; at $c=10^{-4}$ the latter bound exceeds $H(\rho)$, and only \cref{thm:teacher:main} says more.

\begin{proposition}[What no procedure can do]\status{proved; computed}\src{imitation/robust notes-final Props 4.10(a2),(c),(d), 4.11; referee-A R1; checks f2(B), f8, s1(D); referee rb2(B),(C)}\label{prop:imit:limits}
\begin{enumerate}[label=(\alph*),itemsep=1pt,topsep=2pt]
\item \emph{One data law, at most one copy.} If worlds $W_1,\dots,W_N$ as in \cref{def:imit:setting} have one law of the data pairs and persons whose kernels differ pairwise on sets of positive $\bar\mu$-mass, then every procedure $\Psi$ (a measurable map from samples and independent randomness to programs) has $\sum_i\Prob_{W_i}[\Psi(h_n)\in\mathcal E(\nu_i)]\le1$ at every $n$. In the feature case with $|\Sigma|\ge2$, one such world has the person $P^*$, who responds to the cell.
\item \emph{A person entangled with its coupling.} Let $X=W$, $\mu(\cdot\mid w)=\delta_w$, $Y=\Sigma=\{0,1\}$, $\rho=\Bern(s_0/2^m)$ with $1\le s_0<2^{m-1}$, and let the person $H_f$ answer $f(w)\oplus b$ for a table $f:W\to\{0,1\}$. Draw $f$ uniformly and $\beta$ uniformly from its type class, independently. Then $\Prob[\delta_\beta>j]<2^{-j}$, yet every procedure copies the person with probability at most $2^{-L_\beta}$, at every $n$.
\item \emph{Cheap couplings.} In the feature case, with $|\Sigma|\ge2$ and the values of $\rho$ distinct, let $\pi_{\rm u}$, shorter than $\ell_0$, compute the uniform kernel on $Y$. If $H(\rho)/(K(\beta\mid\pi_H)+\kappa)\ge(\log|Y|-H(\rho))/(\ell_0-\ell(\pi_{\rm u}))$, then at every $c>0$ some program that is not a copy has $J_c$ at most that of every copy.
\end{enumerate}
\end{proposition}

\emph{Proof idea.} (a) The output has one law in all the worlds, and the sets of copies are disjoint. (b) The data depend on the world only through $f\oplus\ind_\beta$, which $\binom{2^m}{s_0}$ equally likely worlds with pairwise different persons share. (c) Compare a copy with the program of \cref{thm:imit:main} that computes $P^*$ and with $\pi_{\rm u}$ (\cref{app:imit:limits}).

So imitation data identify behaviour, not mechanism: persons with the same response function and coin law share a data law, however they compute them, and the selected program is a shortest one for the behaviour, which may compute it unlike the person and extrapolates off $\supp\bar\mu$ as that program does. Nor do they identify the split of $P^*$ into a person and an environment: every procedure adopts a convention, and the $c$-procedure adopts it by rate, as steeper selection does for error rules (\cref{prop:limits:separable}). An incompressible coupling is not enough: in (b) the person is as complex as its coupling and responds to the same cells, and a hypothesis such as $\ell_0\ll2^m$, or the left condition that list decoding gives in \cref{thm:teacher:main}, is needed. (c) is a form of rate blindness (\cref{prop:limits:rateblind}): if the person's apparent randomness has a short description, such as a pseudo-random generator with a short seed, no $c$ returns the person with fresh coins; and if that description is the person's own deterministic state, the program that models it is arguably the more faithful copy, so the target $\nu_H$ is itself a modelling choice. In the feature case the same comparison shows that a window for copies needs $2^m>(\ell_0-\ell(\pi_{\rm u}))/(\log|Y|-H(\rho))-(a_\rho+1+\kappa)/H(\rho)$ cells (\cref{app:imit:limits}): with the machine constant set to $0$ (an illustration), a person of $10^6$ bits with $H(\rho)=\frac12$ and $|Y|=2$ needs more than $1.99\cdot10^6$ incompressible cells.

\begin{question}[imitation]\status{open}\src{imitation/robust notes-final \S8 items 1, 2, 4, 6, 7}\label{q:imit:open}
(a) Is there a near-imitation theorem when the coin combines $\beta(w)$ with a private residual coin, or (b) when the cells have unequal masses? (c) Does \cref{thm:imit:main} hold for $cn$-MAPs at each finite $n$ with probability $1-\delta$, for every response function? (d) Is a coupling that is short to describe but slow to compute private randomness for a time-bounded class? (e) Can data other than imitation data, from several environments or from interventions, identify which structure belongs to the person?
\end{question}
